\chapter{סיכום}\label{ch:summary}

תמונות שקריות שנוצרו בבינה מלאכותית כדי לייצר טראפיק אינן בעיה טכנית בלבד וגם לא בעיה אתית בלבד. הן תוצאה של תמריץ כלכלי שפוגש טכנולוגיה זולה, בהיעדר נוהל עבודה. \en{VeriLens} מטפלת בשלושתם: שלוש שכבות אימות מייצרות ראיות, העורך האנושי מחליט ומתעד, התקן \en{TRUTH-IMG} קובע שטראפיק אינו נימוק לפרסום, וממשק ״ציון השלמות״ הופך אמינות לנכס כלכלי.

בדוח זה הוגדרו \num{6} יעדים מדידים, \num{34} דרישות, \num{12} תרחישי בדיקה, \num{10} תוצרים עם קריטריון קבלה לכל אחד, לוח זמנים של \num{18} חודשים ותקציב של \num{9.6} מיליון ₪. הדוח מוגש כחוזה: עם אישורו, כל אלה מחייבים, וכל שינוי ייעשה בבקשת שינוי בלבד.

\section{סוגיות פתוחות לסגירה לפני חודש \num{1}}

לפני שמתחילים נשארו שש סוגיות שעוד לא סגרנו. טבלה~\ref{tab:actions} מפרטת אותן.

\begin{table}[!htb]\centering
\caption{רשימת פעולות פתוחות}\label{tab:actions}
\small
\begin{tabularx}{\linewidth}{@{}l Y l l@{}}
\toprule
\textbf{מזהה} & \textbf{פעולה} & \textbf{אחראי} & \textbf{עד} \\
\midrule
\req{OP-01} & \num{15} עד \num{20} שיחות עם עורכים, בודקי עובדות ומנהלי תביעות בביטוח, לבדיקת הלקוח המשלם (טבלה~\ref{tab:altmarkets}) & גיא ופז & חודש \num{1} \\
\req{OP-02} & תקופת ניסיון אצל \en{ImageWhisperer}, ומיפוי מדויק של מה שחסר בשירות & גיא ופז & חודש \num{1} \\
\req{OP-03} & אישור מאגר \en{VL-Wild} כאמת מידה (\req{A-06}) & מרצה & חודש \num{1} \\
\req{OP-04} & מכתבי כוונות משלוש מערכות חדשות לפחות (\req{A-01}) & גיא ופז & חודש \num{3} \\
\req{OP-05} & חוות דעת משפטית על רישוי צילומים לאימון & ייעוץ משפטי & חודש \num{2} \\
\req{OP-06} & פנייה למועצת העיתונות או לגוף אתיקה אחר, להקמת ועדת ההסמכה העצמאית (סעיף~\ref{sec:enforce}) & גיא ופז & חודש \num{4} \\
\bottomrule
\end{tabularx}
\end{table}

\begin{keybox}[מבוקש]
אישור המרצה לדוח המכין, והרשאה להתחיל בשלב \num{1} (מאגרי נתונים) ובכתיבת טיוטת התקן.
\end{keybox}
